\documentclass[10pt,a4paper]{amsart}
\usepackage[margin=19mm]{geometry}
\usepackage{amsmath,amssymb,mathtools}
\usepackage[scr=rsfs,cal=euler]{mathalfa}
\usepackage{xfrac}
\usepackage[unicode,hidelinks,bookmarks=false]{hyperref}
\newcommand{\ie}{i.e.\ }
\newcommand{\cf}{cf.\ }
\newcommand{\resp}{respectively\ }
\newcommand{\Subsection}{\subsection}
\providecommand{\nobreakdash}{\nobreak\hbox{-}}
\setlength{\emergencystretch}{2em}
\multlinegap0pt
\usepackage{amssymb}
\usepackage[shortlabels]{enumitem}
\newtheorem{lemma}{Lemma}
\usepackage{cleveref}
\crefname{lemma}{Lemma}{Lemmas}
\newcommand{\A}{\mathbb{A}}
\newcommand{\CC}{\mathbb{C}}
\newcommand{\GL}{\mathrm{GL}}
\renewcommand{\Im}{\mathrm{Im}}
\newcommand{\Q}{\mathbb{Q}}
\newcommand{\R}{\mathbb{R}}
\newcommand{\diag}{\mathrm{diag}}
\title{Tatuzawa's theorem for Rankin--Selberg $L$-functions}
\author{Gergely Harcos, Jesse Thorner}
\date{Source: arXiv:2508.10844v3 (CC BY 4.0)}
\begin{document}
\maketitle

\noindent\emph{Source and licence.} Tatuzawa's theorem for Rankin--Selberg $L$-functions. Version arXiv:2508.10844v3 (CC BY 4.0), distributed by arXiv under the Creative Commons Attribution 4.0 International licence, http://creativecommons.org/licenses/by/4.0/, as recorded in the arXiv licence metadata for this exact version and shown on the abstract page https://arxiv.org/abs/2508.10844v3. Full licence notice: \texttt{licenses/arxiv-2508.10844-CC-BY-4.0.txt}.\par\smallskip

\noindent\textit{Editorial scope.} Counted unit: the article's lemma tpiformula (source0.tex lines 332--346). The article's complete original proof of it is delivered with it (source0.tex lines 348--362, one \texttt{proof} environment of 15 lines). No mathematics is written, repaired or re-derived: the statement and the proof are the article's own lines, in the article's own notation, and no retained line is substituted or re-flowed.  Retained uncounted as the source-local context the proof needs: lines 252--255; lines 269--271; lines 297--306.  Nothing was omitted from any retained span, so the package carries no omission record.  No retained line contains a citation, so the wrapper carries no bibliography and no bibliography entry.  No retained line contains a figure environment, an image-inclusion command or drawing code, so no image is delivered and the package declares no asset dependency.  Editorial formatting only, in addition: an AMS \texttt{amsart} preamble that restates the article's own declarations and macros used by the retained lines, namely the environments \texttt{lemma} on separate counters, together with the standard packages those lines need.  The retained environments print in this document as Lemma 1 in order of appearance, whereas the article prints the same items with the numbering of its own full document.  The complete native source of the article as distributed in the arXiv e-print for this exact version is bundled unchanged as \texttt{sources/183283/source0.tex}.
\begin{equation}
\label{eqn:pidecomp}
\pi=\pi^*\otimes|\cdot|^{it_{\pi}},\qquad \pi^*\in\mathfrak{F}_n^*,\quad t_{\pi}\in\R.
\end{equation}

\begin{equation}\label{eq:Lpiv}
L(s,\pi_{v})=\prod_{j=1}^n \Gamma_v(s+\mu_{j,\pi}(v)).
\end{equation}

\begin{lemma}\label{lem:rhochar}
Let $K\in\{\R,\CC\}$, and let $\rho$ be a generic irreducible admissible representation of $\GL_n(K)$ with central character $\omega_\rho$. There exist complex numbers $\mu_1,\dotsc,\mu_n\in\CC$ such that
\begin{equation}\label{eq:Lrho}
L(s,\rho)=\prod_{j=1}^n\Gamma_K(s+\mu_j)
\end{equation}
and
\begin{equation}\label{eq:wrho}
\frac{\omega_\rho(a)}{|\omega_\rho(a)|}=\prod_{j=1}^n a^{i[K:\R]\Im(\mu_j)},\qquad a>0.
\end{equation}
\end{lemma}

\begin{lemma}
\label{tpiformula}
Let $\pi\in\mathfrak{F}_n$ with archimedean Langlands parameters $\mu_{j,\pi}(v)$ as in \eqref{eq:Lpiv}.  If
\begin{equation}\label{eq:mpi}
m_{\pi}(v)=[F_v:\R]\sum_{j=1}^n \Im(\mu_{j,\pi}(v)),
\end{equation}
then
\begin{equation}\label{eq:tpi}
n[F:\Q]t_\pi=\sum_{v\mid\infty}m_{\pi}(v).
\end{equation}
In particular,
\begin{equation}\label{eq:tpibound}
|t_\pi|\leq\max_{\substack{v\mid\infty\\1\leq j\leq n}}|\Im(\mu_{j,\pi}(v)|.
\end{equation}
\end{lemma}

\begin{proof}
Since the central character $\omega_\pi$ is unitary, it follows from \eqref{eq:Lpiv}, \eqref{eq:mpi}, and \cref{lem:rhochar} that
\[
\omega_{\pi_v}(a)=a^{im_{\pi}(v)},\qquad v\mid\infty,\quad a>0.
\]
Now consider the archimedean central elements
\[
D(a)=\prod_{v\mid\infty}\diag(a,\dotsc,a)\in\GL_n(\A),\qquad a>0.
\]
Observe that $\det(D(a))$ has idele norm $a^{n[F:\Q]}$, and recall the definition of $t_\pi$ implicit in \eqref{eqn:pidecomp}. Calculating the action of $\pi(D(a))$ in two ways, we see that
\[
(a^{n[F:\Q]})^{it_\pi}=\omega_\pi(a)=\prod_{v\mid\infty}\omega_{\pi_v}(a)=\prod_{v\mid\infty}a^{im_{\pi}(v)},\qquad a>0.
\]
Comparing the two sides, the identity \eqref{eq:tpi} and the bound \eqref{eq:tpibound} follow readily.
\end{proof}

\end{document}
